\section{语义进展表示与历史条件纠偏}
\label{sec:method}
\subsection{总体信息流与训练边界}
方法以ACT作为慢策略，输出基础动作块；以轻量GRU为快分支，输出相同目标槽位的
局部修正。两者承担不同职责：慢策略组织连贯运动并编码任务几何，快分支利用
更近的观测和历史调整当前计划。动作空间残差的通用思想已有研究\cite{zeng2020}，
动作块实时修正也已有专门工作\cite{a2c2}。本文研究语义监督、条件历史与时间对齐的
组合如何影响遮挡接触任务中的动作块执行，而非重新提出一般残差控制结构。

训练分三步：先构建标签并训练分割器；再冻结分割器训练语义ACT及几何query；
最后冻结整个基础策略，只训练快分支。在线不运行SAM2或读取数据集mask。
三个阶段共享预先确定的episode划分，任何作为全系统未见测试的episode不得用于
分割训练、提示器训练、策略训练、阈值选择或checkpoint选择。
若已有checkpoint不满足统一划分，其结果只能按实际已见条件报告。

\subsection{把伪标签噪声控制前移}
为降低式\eqref{eq:noise}中的标签噪声，首先将类别说明与代表帧转换为AI生成的
稀疏框/点JSON提示，渲染后人工审核。SAM2\cite{ravi2024sam2}传播少量种子视频；
对目标初现、接触及再遮挡片段补充提示并局部重跑。审核标签转为检测框后训练
分视角YOLO提示器，用于扩大标注范围，再进入诊断与修复循环。

诊断包括邻帧IoU、面积突变、质心跳变、连通域异常，以及可用时的正反传播一致性。
各项按类别阈值归一化为分数，连同预期可见性、缺失项和人工例外独立保存。
低可信片段进入自动重试或少量人工复核；真实暴露事件不能仅因面积突变被删除。
分数用于定位风险而非宣称已校准的不确定概率；一致但错误的类别仍需独立审核。
双向结果使用未来帧，严格限制在离线标注中。

标签质量影响三个不同环节：离线修复、提示器样本构建、以及下述几何辅助监督。
当前独立U-Net训练采用类别加权CE与配置的前景Dice，不默认把每帧质量分数
自动乘入分割loss。论文需分别记录各环节实际启用的机制，避免将整个流程笼统称为
概率不确定性端到端学习。人工时间、提示修改次数和残余标签错误应共同计入成本收益。

\subsection{共享视觉编码的语义动作块策略}
分别为front和side训练轻量U-Net，保留对应视角类别数，不强制类别数相同。
当前任务front含背景共7类，side含背景共5类；同名类别使用一致调色板。
完整视野缩放到 $480\times270$ 后分割，策略训练和运行中均冻结分割器。
对概率 $P_t^v(c,u)$，语义图定义为
\begin{equation}
 S_t^v(u)=\sum_{c=1}^{C_v}P_t^v(c,u)\,\bm p_c,
 \quad \bm p_c\in[0,1]^3.
 \label{eq:palette}
\end{equation}
这是概率加权颜色，不是硬argmax；多类别到三通道的映射一般不可逆。
每个视角的RGB和语义图分别经过同一个ResNet18及投影，作为独立视觉token输入
ACT encoder，保留原ACT位置编码，不增加camera或modality embedding。
单视角包含两路视觉输入，双视角包含四路。保留RGB是为了避免语义瓶颈
丢弃姿态、机械臂外观及其他没有标入类别的控制信息。

本结构是对几何归纳偏置的实用实现，不声称共享编码器最优。
同参数或同token预算的重复RGB对照应排除仅由额外计算带来的收益。
当前不进行跨视角三维配准，也不将side分割变换为补全后的front图。

\subsection{用几何监督定义query职责}
在encoder中加入三个可学习初始向量 $e_i\in\R^{512}$。
它们与视觉、关节及ACT原有token共同进行自注意力，产生观察相关表示 $z_i$。
线性头 $d_i(z_i)$ 分别预测以下front图像几何：
\begin{align}
 G_1&=|M_O|/(W H_I),\nonumber\\
 G_2&=\norm{c_O-c_R}/\sqrt{(W-1)^2+(H_I-1)^2},\nonumber\\
 G_3&=\bar c_O-\bar e_{T,L}\in\R^2.
 \label{eq:targets}
\end{align}
$M_O$ 为目标可见mask，$c_O,c_R$ 为目标和区域质心；横纵归一化坐标用上横线表示。
将tool mask在归一化平面上做PCA主轴拟合，取投影极值的两个端点，
选择横坐标较小者 $\bar e_{T,L}$。该定义保留目标与工具端点的二维方向关系，
但图像左端点不必总是实际接触端，需在对应装置与视角下验证。

目标不可见时仍监督面积为零；组成关系的任一mask不可见时不监督相应距离或向量。
令 $v_{bi}$ 表示该组可定义，$q_{bi}$ 为相关类别质量的最小值，
$\omega_{bic}=v_{bi}\min(q_{bi}/0.95,1)$ 为复制到各坐标 $c$ 的权重，则
\begin{equation}
 L_G=\frac13\sum_{i=1}^3
 \frac{\sum_{b,c}\omega_{bic}\ell_{0.01}(\widehat G_{bic}-\widetilde G_{bic})}
 {\max(1,\sum_{b,c}\omega_{bic})}.
 \label{eq:qloss}
\end{equation}
此处 $\widehat G_{bi}=d_i(z_{bi})$，三组维度为1、1、2；先按有效坐标权重归一化，
再对三个组取平均，与当前实现一致。
由于分母同样加权，所有样本同倍降权可能被抵消，不能称其为整体loss的绝对衰减。
ACT学习目标为
\begin{equation}
 L_{\mathrm{slow}}=L_{\mathrm{action,L1}}
       +\lambda_{\mathrm{KL}}L_{\mathrm{KL}}+\lambda_G L_G,
\end{equation}
当前配置 $\lambda_G=1$。动作decoder读取完整encoder上下文，包括query隐藏状态。
几何loss更新ACT表示但不更新冻结U-Net。

该设计对应式\eqref{eq:geobound}的几何读出路径；要得到严格比较，
还需式\eqref{eq:geoguarantee}中的基线信息损失与实际读出误差条件。
它不是监督attention map，也不是把预测几何四个数直接作为唯一动作输入。
辅助头是否准确与动作是否依赖query必须分别检验。
当前模型不额外学习阶段分类、切换或停止头；可见性、分离和区域包含等事件
用于独立评价其任务进展行为，而不是添加未经实现的控制规则。

\subsection{最新观察与短期历史的增量反馈}
基础策略冻结后，每次快更新只处理新到达的双视角RGB，完整视野缩放到
$320\times180$，复用冻结视觉骨干并池化为每视角 $2\times3\times512$ 特征。
不重复编码窗口内旧图像。控制上下文 $c_t$ 包括实际follower关节、因果估计速度、
上一发送指令、对应基础未来动作及相对当前关节的差、槽位有效性、计划年龄、
观察间隔与计划切换信息。显式提供基础目标动作，使快分支预测的是该计划的误差。

以三个分支进行模块化研究：A使用普通RGB ACT及最新RGB；B使用语义ACT，并在
每次快更新运行双U-Net得到最新语义描述；C使用QToken ACT，将\emph{当前生效慢计划}
的三个query隐藏状态作为语义上下文，并结合最新RGB。
B额外描述包含概率池化、面积、质心、可见性及关系有效性；C的query并非每次快更新
重新计算，也不是最新语义观测。数值有效标志不能充当分割可信度。

三者使用共同的历史编码形式。对最近 $N_h=8$ 条记录
\begin{equation}
 h_i=\operatorname{GRU}_\theta(\rho_\theta(f_i,c_i),h_{i-1}),
 \quad h_{\mathrm{start}-1}=0.
 \label{eq:history}
\end{equation}
视觉及输入投影为128维，GRU为单层hidden128。
每次从窗口起点重算GRU，不跨调用无限延续隐藏状态；5Hz时8条记录覆盖约1.4秒。
reset或超过0.3秒的记录断流清空窗口。历史编码旨在降低式\eqref{eq:unified}的
有效编码误差，但其实际收益须与同输入最新记录MLP、历史置乱和窗口长度对照区分。

\subsection{有界动作残差与监督}
快头输出 $H$ 个动作修正，在机器人原生命令单位中定义
\begin{equation}
 \delta a_{t,j}=\bm\beta\odot\tanh([W_rh_t+c_r]_j),
 \quad \widehat a_{t,j}=b^{(k)}_{t+d+j}+\delta a_{t,j}.
 \label{eq:action}
\end{equation}
输出层零初始化，初始等于基础策略。每维 $\beta_i$ 由训练基础误差绝对值P95设定，
当前裁剪到 $[0.05,5]$ 原生命令单位；换平台必须核对单位和尺度，不能直接称为角度。
限幅约束局部修改范围，不保证动力学安全。

令 $m_{tj}$ 为有效监督槽位，动作维度为 $n_a$。定义标准差缩放误差
$e_{tji}=(\widehat a_{tji}-a^{\mathrm{demo}}_{t+d+j,i})/\sigma_i$、
$r_{tji}=\delta a_{tji}/\sigma_i$，学习目标为
\begin{equation}
 L_R=\frac{\sum_{t,j,i}m_{tj}
 [\ell_{0.1}(e_{tji})+\lambda_R r_{tji}^2]}
 {n_a\max(1,\sum_{t,j}m_{tj})}.
 \label{eq:resloss}
\end{equation}
各项按动作维度归一化，当前 $\lambda_R=10^{-3}$。
训练数据使用leader示范指令作为目标，follower关节作为观测，不混用此前delta实验。
只训练新增投影、GRU和残差头；特征可按episode离线缓存，尾帧屏蔽无效目标。
上一指令训练时取示范，实机时为实际发送值，这一分布差异需在闭环中评价。
AdamW学习率 $10^{-4}$、batch128、最多50轮，验证8轮无改善早停；
按验证标准差缩放MAE选checkpoint，平方风险另作机制诊断。

\subsection{执行时间对齐与局部重叠}
当前实例使用30Hz控制，慢策略输出60步、每30步更新基础计划；快分支每6帧更新，
延迟偏移 $d=3$，输出 $H=9$。例如第6帧观察对应第9至17帧修正，
第12帧观察对应第15至23帧修正。相邻轮次重叠3帧；及时应用首槽的端到端预算为100ms。

所有结果携带plan ID、观察时间和目标槽位。新结果仅替换同计划同槽位的未执行残差，
始终加到原基础动作上，不向已经修正的动作累加。迟到按式\eqref{eq:delay}丢弃前缀；
新计划生效清空旧残差，重复、非有限和失效观察不更新反馈。
当前接受结果的最大观察年龄0.15秒、残差使用年龄上限0.4秒。
因此实际接受规则可能比纯覆盖界更严格，必须计入实测可用覆盖率。
所有参与结果接受的时序状态均属于运行记录；分析中应纳入相应可用信息集。

残差无效而基础槽位有效时回退原计划；两者都无效时由硬件层执行已验证的安全处置，
不能无限复用最后一个动作。端到端预算包括图像传输、预处理、编码、排队及命令发布，
不能以GRU单次GPU耗时替代。本文方法规格须与每次实机部署版本核对；离线shadow
通过不自动等价于真实硬件实时验证。
